\documentclass[10pt,a4paper]{article}
\usepackage[utf8]{inputenc}
\usepackage[T1]{fontenc}
\usepackage[spanish,es-noshorthands]{babel}
\usepackage[margin=1.8cm]{geometry}
\usepackage{booktabs,array,enumitem,amsmath,amssymb,xcolor,tcolorbox,icomma,multicol}
\usepackage[hidelinks]{hyperref}
\pagestyle{plain}
\setlength{\parindent}{0pt}\setlength{\parskip}{4pt}
\newtcolorbox{ejemplar}{colback=green!5,colframe=green!40!black,title=Respuesta ejemplar,fonttitle=\bfseries\small,boxrule=0.6pt,left=4pt,right=4pt,top=2pt,bottom=2pt}
\newtcolorbox{trampa}{colback=red!5,colframe=red!50!black,title=Trampas,fonttitle=\bfseries\small,boxrule=0.6pt,left=4pt,right=4pt,top=2pt,bottom=2pt}
\newtcolorbox{enun}[1]{colback=blue!4,colframe=blue!45!black,title={#1},fonttitle=\bfseries\small,boxrule=0.6pt,left=4pt,right=4pt,top=2pt,bottom=2pt}
\newtcolorbox{gancho}{colback=yellow!8,colframe=orange!70!black,boxrule=0.6pt,left=5pt,right=5pt,top=3pt,bottom=3pt}
\newcommand{\V}{\textsuperscript{\textcolor{green!45!black}{[V]}}}
\newcommand{\Pre}{\textsuperscript{\textcolor{orange!80!black}{[P]}}}
\newcommand{\inv}{\textsuperscript{\textcolor{red!60!black}{[inv]}}}
\newcommand{\calc}{\textsuperscript{\textcolor{blue!60!black}{[cálc]}}}
\newcommand{\acto}[1]{\vspace{6pt}{\large\bfseries #1}\par\vspace{2pt}\hrule\vspace{4pt}}
\newcommand{\Ph}{\hat P}
\begin{document}
\renewcommand{\%}{\ifmmode\mbox{\char37}\else\char37\fi}

\begin{center}
{\LARGE\bfseries Borrón y cuenta nueva: la Ley Dicom de 2012}\\[3pt]
{\large Logit, clasificación, \emph{scoring} y diferencias en diferencias sobre un caso real del crédito chileno}\\[3pt]
\small Ficha de estudio de la parte teórica · los 22 ítems seleccionados (núcleo 12 + refuerzo 10) · mar 3-11-2026 · PEP 1: vie 6-11
\end{center}

\textbf{Cómo se usa.} Cada ítem del banco de teoría aparece \emph{reescrito sobre el caso} (recuadro azul), con su respuesta ejemplar (verde) y sus trampas (rojo). Primera pasada: lee el caso de cada acto. Segunda pasada: tapa lo verde, contesta en voz alta (V/F + palabra clave) y compara. La última página es la tabla de recuerdo activo. En la prueba los enunciados hablarán de Spector, deserción, PAES o SLEP: \textbf{el razonamiento es el mismo}, solo cambian los nombres.

\textbf{Marcas de las cifras.} \V\ leída en la fuente primaria · \Pre\ prensa o terceros · \calc\ cálculo propio desde cifras \V · \inv\ inventada para practicar (la fuente no publica ese dato), coherente con lo real. Nada inventado se presenta como real.

\begin{gancho}
\textbf{Enero--febrero de 2012.} El Congreso aprueba la Ley 20.575, conocida como «Ley Dicom»: los burós de crédito dejan de informar las deudas impagas que se hicieron exigibles antes del 31-12-2011, si el total de la persona (sin intereses ni reajustes) es menor a \textbf{\$2,5 millones} (unos US\$5.000)\V\Pre. Además, la información comercial no puede usarse para seleccionar personal, admitir en colegios ni condicionar atenciones de salud\Pre. Se justificó como alivio tras el terremoto de 2010, y ya había existido un borrón parecido en 2002\V.

\textbf{Lo que pasó.} Se borraron los registros de \textbf{2,8 millones de personas}, el \textbf{21\%} de los 13 millones de chilenos mayores de 15 años y el \textbf{84\%} de los deudores bancarios en mora\V. Los bancos siguieron viendo la mora \emph{bancaria} (el regulador, la SBIF, no borró nada), pero ya no la mora en el comercio\V. La tasa mediana de los créditos de consumo pequeños subió \textbf{5,3 pp} (desde 26\%)\V. Tres años después, el número de personas con mora informada casi había vuelto al nivel previo\V.

\textbf{La pregunta del economista:} ¿la ley amplió el acceso al crédito? Liberman, Neilson, Opazo y Zimmerman (NBER w25097) responden con los datos de \emph{todos} los deudores bancarios de Chile (Sinacofi, 2009--2014; 23 millones de observaciones de 5,6 millones de personas)\V. Para eso usan las cuatro herramientas de la PEP: un modelo de \textbf{probabilidad de impago} (Acto 1), su \textbf{capacidad de clasificar} (Acto 2), la lógica del \textbf{score} (Acto 3) y un \textbf{DiD} (Acto 4). Respuesta corta: el crédito total \textbf{cayó 3,5\%}\V.
\end{gancho}

% =====================================================================
\acto{Acto 1. ¿Quién no paga? Logit y MPL del impago bancario}

\textbf{El caso.} Para cada deudor, los autores predicen el \textbf{impago bancario en los 6 meses siguientes} (mora de 90 días o más), con y sin el registro borrado\V. Antes del borrón, la tasa de impago era \textbf{0,029} en el mercado de «bajo costo» (casi todos sin registro en Dicom) y \textbf{0,069} en el de «alto costo» (casi todos con registro)\V. Con una sola dummy, $REG=1$ si la persona tenía mora informada, el logit reproduce exactamente las dos tasas, así que sus coeficientes salen a mano\calc:
\[
\hat z=-3{,}511+0{,}909\,REG,\qquad \Ph=\frac{1}{1+e^{-\hat z}}\quad (y=1 \text{ si no paga, es decir, «malo»}).
\]
Comprueba: $REG=0\Rightarrow \Ph=1/(1+e^{3,511})=0{,}029$; \ $REG=1\Rightarrow \hat z=-2{,}602$, $\Ph=0{,}069$.

\begin{enun}{T14 (núcleo) · signo con $y=1$ malo}
V/F. En el probit del Banco Central de Chile sobre impago de los hogares (EFH 2007), $y=1$ si el hogar no paga, y el coeficiente de la carga financiera (RCIM, hasta 42,8) es positivo. Eso indica que a mayor carga financiera, menor riesgo.
\end{enun}
\begin{ejemplar}
\textbf{F.} Con $y=1=$ malo, un $\hat\beta>0$ significa \textbf{más} probabilidad del evento malo: más carga financiera, \textbf{más riesgo}. Martínez, Cifuentes, Madeira y Poblete-Cazenave (2013) obtienen $\hat\beta>0$ para la carga financiera bajo 42,8 (sobre ese umbral la pendiente deja de ser significativa) y para la probabilidad de despido del jefe de hogar ($8{,}116$); el ingreso tiene $\hat\beta=-0{,}216$: más ingreso, menos riesgo\V. Lo mismo en nuestro logit: $\hat\beta_{REG}=+0{,}909$, tener mora informada sube el riesgo. En el puntaje (score) es al revés: más puntos, menos riesgo.
\end{ejemplar}

\begin{enun}{T6 (núcleo) · $\hat\beta$ no está en puntos porcentuales}
V/F. $\hat\beta_{REG}=0{,}909$ significa que tener mora en Dicom aumenta la probabilidad de impago en 0,91 puntos porcentuales.
\end{enun}
\begin{ejemplar}
\textbf{F.} Un coeficiente logit mueve el \emph{índice} $z$ (log-odds), no la probabilidad. Lo interpretable es $e^{0,909}=2{,}48$: tener registro \textbf{multiplica los odds de impago por 2,48}\calc. En probabilidad, como $REG$ es dummy, el efecto es la resta de dos $\Ph$: $0{,}069-0{,}029=\mathbf{+4{,}0}$ \textbf{pp}, no 0,91.
\end{ejemplar}
\begin{trampa}
\begin{itemize}[leftmargin=*,itemsep=0pt,topsep=0pt]
\item Para una dummy, el efecto es $\Ph(1)-\Ph(0)$, no $g(z)\hat\beta$.
\item «Odds ×2,48» no es «probabilidad ×2,48»: la probabilidad se multiplica por $0{,}069/0{,}029=2{,}38$\calc. Se parecen solo porque $P$ es chica.
\end{itemize}
\end{trampa}

\begin{enun}{T13 (refuerzo) · pp frente a \%}
V/F. Tras el borrón, la tasa mediana de los créditos pequeños pasó de 26\% a 31,3\%: subió 5,3\%.
\end{enun}
\begin{ejemplar}
\textbf{F.} Subió \textbf{5,3 puntos porcentuales}; en términos relativos, $5{,}3/26=\mathbf{20\%}$, que es justo lo que dice el paper\V\calc. Lo mismo con el impago: de 2,9\% a 6,9\% son \textbf{4,0 pp}, o $+138\%$ relativo\calc. «pp» es resta; «\%» es cambio relativo. Es el error de «de 0,04 a 0,06 sube 2\%» (son 2 pp, y +50\%).
\end{ejemplar}
\begin{trampa}
\begin{itemize}[leftmargin=*,itemsep=0pt,topsep=0pt]
\item Con logaritmos, una diferencia de $-0{,}29$ en $\ln \Ph$ es $e^{-0,29}-1=-25\%$\calc, aunque el paper lo redondee a «cae 29\%»\V. La aproximación $\Delta\ln\approx\%$ solo vale para cambios chicos.
\end{itemize}
\end{trampa}

\begin{enun}{T1 (refuerzo) · heterocedasticidad del MPL}
V/F. Si un banco estima un MPL, $\Ph=0{,}029+0{,}040\,REG$, como $\mathrm{Var}(u\mid x)=P(1-P)$ cambia entre grupos, los $\hat\beta$ de MCO quedan sesgados.
\end{enun}
\begin{ejemplar}
\textbf{F.} La varianza sí cambia: $0{,}029\cdot0{,}971=0{,}028$ sin registro y $0{,}069\cdot0{,}931=0{,}064$ con registro, \textbf{2,3 veces} más\calc. Pero la heterocedasticidad \textbf{no sesga} a MCO: le quita \textbf{eficiencia} e invalida los errores estándar habituales. Se corrige con errores robustos (o MCP con $\hat y(1-\hat y)$). Con una sola dummy, el MPL reproduce exactamente las dos tasas.
\end{ejemplar}

\begin{enun}{T22 (núcleo) · mismo $\hat\beta$, distinto efecto}
Comente. «En el probit del Banco Central, un punto porcentual más de probabilidad de despido tiene el mismo efecto en un hogar con $\Ph=0{,}10$ que en uno con $\Ph=0{,}50$, porque $\hat\beta=8{,}116$ es uno solo».
\end{enun}
\begin{ejemplar}
\textbf{Falso.} El efecto marginal es $g(z)\hat\beta$, con $g$ la densidad (normal en el probit; $P(1-P)$ en el logit). Con $\Ph=0{,}10$: $z=-1{,}28$, $\phi(z)=0{,}175$, y $0{,}175\cdot8{,}116\cdot0{,}01=\mathbf{+1{,}4}$ \textbf{pp}. Con $\Ph=0{,}50$: $\phi(0)=0{,}399$, y $\mathbf{+3{,}2}$ \textbf{pp}, más del doble\calc. El hogar «en el margen» ($P\approx0{,}5$) es el que más responde a un shock laboral. Por eso el Banco Central mide el riesgo por estrato y no con un solo número: con la misma carga financiera, el estrato de menores ingresos tiene más probabilidad de impago\V.
\end{ejemplar}

\textbf{Una variable continua\inv.} Para el ítem siguiente se agrega $CF=$ carga financiera (\% del ingreso que va a cuotas): $\hat z=-4{,}0+0{,}91\,REG+0{,}025\,CF$\inv. En el logit, $g(z)=P(1-P)$.

\begin{enun}{T11 (núcleo) · EPeP frente a EPP}
V/F. Con cuatro deudores\inv\ ($REG,CF$) $=(0,10),(0,30),(1,20),(1,60)$, el efecto parcial en el promedio y el efecto parcial promedio de $CF$ coinciden, porque salen del mismo modelo.
\end{enun}
\begin{ejemplar}
\textbf{F.} \textbf{EPeP:} se evalúa $g$ en el deudor promedio, $\bar z=-4{,}0+0{,}91(0{,}5)+0{,}025(30)=-2{,}795$, $\Ph=0{,}058$, $g(\bar z)=0{,}054$, y $0{,}054\cdot0{,}025=\mathbf{+0{,}136}$ pp por punto de carga. \textbf{EPP:} se promedian los cuatro efectos individuales, con $\Ph_i=0{,}023;\ 0{,}037;\ 0{,}070;\ 0{,}169$, $\overline{g}=0{,}066$, y $0{,}066\cdot0{,}025=\mathbf{+0{,}165}$ pp\calc. Como $g$ no es lineal, $g(\bar z)\neq\overline{g(z_i)}$. Además, nadie tiene $REG=0{,}5$: el «deudor promedio» no existe, que es el argumento a favor del EPP.
\end{ejemplar}

% =====================================================================
\acto{Acto 2. ¿Sirve el modelo para separar? Clasificación con tasa base baja}

\textbf{El caso.} En la muestra con crédito nuevo, la tasa de impago bancario a 6 meses es \textbf{5\%}\V: los «malos» son pocos. Los autores comparan tres predictores (bosque aleatorio, \textbf{LASSO logístico} y Bayes ingenuo) por su log-verosimilitud fuera de muestra; el bosque gana y, al sacarle el registro borrado, empeora de $-0{,}278$ a $-0{,}284$\V\ ($-2{,}2\%$)\calc. El paper no publica matriz de confusión ni AUC, así que la siguiente matriz es \textbf{inventada}, con la tasa base real\inv: $N=10.000$ deudores, 500 malos.
\begin{center}\small
\begin{tabular}{lcc|c}\toprule
Corte 0,5: Real $\backslash$ Predicho & Malo & Bueno & Total\\\midrule
Malo ($y=1$) & VP $=40$ & FN $=460$ & 500\\
Bueno ($y=0$) & FP $=20$ & VN $=9.480$ & 9.500\\\bottomrule
\end{tabular}
\end{center}

\begin{enun}{T33 (núcleo) · pareo de las métricas, con estos números}
Calcula las métricas de la matriz y di qué es la fórmula que sobra, FP/(FP + VP).
\end{enun}
\begin{ejemplar}
Exactitud $=(VP+VN)/N=9.520/10.000=\mathbf{95{,}2\%}$ · Tasa de error $=(FP+FN)/N=4{,}8\%$ · Sensibilidad $=VP/(VP+FN)=40/500=\mathbf{8\%}$ · Especificidad $=VN/(VN+FP)=9.480/9.500=\mathbf{99{,}8\%}$ · Precisión $=VP/(VP+FP)=40/60=66{,}7\%$ · F1 $=2PS/(P+S)=14{,}3\%$\calc. La que sobra, $FP/(FP+VP)=33{,}3\%$, es $1-$ precisión (de los que el modelo llama malos, cuántos pagan).
\end{ejemplar}

\begin{enun}{T34 (núcleo) · exactitud frente a tasa base}
Comente. «El modelo de riesgo del banco acierta el 95,2\% de los casos: es muy bueno».
\end{enun}
\begin{ejemplar}
Un «modelo» que dijera \textbf{nadie es malo} acertaría $9.500/10.000=\mathbf{95\%}$\calc: el modelo le gana por \textbf{0,2 pp}. Con clases desbalanceadas la exactitud no prueba nada. Lo que importa al banco es la \textbf{sensibilidad}: detecta solo 40 de los 500 malos (8\%); los otros 460 reciben el crédito y no pagan. Hay que mirar sensibilidad, especificidad y comparar modelos con la ROC y el AUC (o el KS).
\end{ejemplar}
\begin{trampa}
\begin{itemize}[leftmargin=*,itemsep=0pt,topsep=0pt]
\item \textbf{Caso profesional real: recesiones.} La Fed de Nueva York publica cada mes $P(\text{recesión en 12 meses})=\Phi(-0{,}5333-0{,}6330\cdot spread)$, un probit sobre la diferencia de tasas a 10 años y 3 meses, estimado con datos de 1959--2009 (Estrella y Mishkin, 1998)\V\Pre. En esos años, EE.\,UU.\ estuvo en recesión 93 de 612 meses: \textbf{15\%}\calc. «Nunca habrá recesión» acierta el 85\%.
\end{itemize}
\end{trampa}

\begin{enun}{T23 (refuerzo) · el corte 0,5}
V/F. El banco clasifica como malo si $\Ph>0{,}5$; ese corte es indiscutiblemente el mejor.
\end{enun}
\begin{ejemplar}
\textbf{F.} El 0,5 es el criterio elemental, y con 5\% de malos casi nadie lo supera. El mejor corte depende de los \textbf{costos}: un FN es un malo aprobado (pérdida del préstamo); un FP es un bueno rechazado (negocio perdido). Si se baja el corte a la tasa observada, $\bar y=0{,}05$, con la matriz\inv\ VP $=350$, FN $=150$, FP $=2.400$, VN $=7.100$: la sensibilidad sube a \textbf{70\%}, la especificidad baja a \textbf{74,7\%} y la exactitud cae a 74,5\%\calc. Peor exactitud, modelo más útil. En el modelo de la Fed, el corte 0,5 equivale a $spread<-0{,}84$; un corte de 0,3 equivale a $spread<0$, es decir, a la curva invertida\calc.
\end{ejemplar}
\begin{trampa}
\begin{itemize}[leftmargin=*,itemsep=0pt,topsep=0pt]
\item \textbf{Un falso positivo real:} en julio de 2023 el modelo de la Fed daba \textbf{65--66\%} de probabilidad de recesión\V; el NBER no ha fechado ninguna recesión posterior a la de 2020\V. La curva invertida también dio falsas alarmas a fines de los 60 y de los 90\V.
\item En los sistemas de alerta temprana de crisis (Kaminsky y Reinhart, 1999), el umbral de cada indicador se elige minimizando la razón ruido/señal $=\frac{FP/(FP+VN)}{VP/(VP+FN)}=\frac{1-\text{especificidad}}{\text{sensibilidad}}$: otra forma de elegir el corte por costos.
\end{itemize}
\end{trampa}

\begin{enun}{T27 (refuerzo) · ejes de la ROC}
V/F. La ROC del modelo de riesgo grafica la sensibilidad contra la especificidad para cada corte.
\end{enun}
\begin{ejemplar}
\textbf{F.} Grafica la sensibilidad (eje $y$) contra \textbf{$1-$especificidad} (eje $x$, tasa de falsos positivos) para cada corte. Corte 0,5: punto $(0{,}002;\ 0{,}08)$; corte 0,05: punto $(0{,}253;\ 0{,}70)$\calc. El AUC resume todos los cortes (0,5 sería tirar una moneda). Lectura del caso: la Ley Dicom no cambió el corte; \textbf{bajó la curva}, porque quitó una variable informativa. Eso es lo que miden los autores con la log-verosimilitud.
\end{ejemplar}

% =====================================================================
\acto{Acto 3. El score con y sin Dicom: WOE e IV}

\textbf{El caso.} La variable «monto en mora informado (dic.\ 2011)» tiene cuatro atributos. Las proporciones de deudores son reales: 63\% sin registro, 31\% con mora bajo \$2,5 millones (borrable) y 6\% sobre ese monto\V; el corte del 31\% en dos tramos y las tasas de impago son \textbf{inventados}\inv, coherentes con la tasa base de 5\%. Convención de la clase: $WOE=\ln(\%B/\%M)\times100$, buenos arriba; IV sin $\times100$\calc:
\begin{center}\small
\renewcommand{\arraystretch}{1.15}
\begin{tabular}{lrrrrrr}\toprule
Atributo & $N$ & Impago & \%B & \%M & WOE & IV parcial\\\midrule
Sin registro & 63.000 & 3,0\% & 64,29 & 38,18 & $+52{,}1$ & 0,136\\
Mora hasta \$200 mil & 18.000 & 6,5\% & 17,71 & 23,64 & $-28{,}9$ & 0,017\\
Mora \$200 mil a \$2,5 M & 13.000 & 9,0\% & 12,45 & 23,64 & $-64{,}1$ & 0,072\\
Mora \$2,5 M o más & 6.000 & 12,0\% & 5,55 & 14,55 & $-96{,}3$ & 0,087\\\midrule
Total & 100.000 & 4,95\% & 100 & 100 & & \textbf{0,31}\\\bottomrule
\end{tabular}
\end{center}
Comprueba la fila 1: buenos $=63.000\cdot0{,}97=61.110$, el 64,29\% de los 95.050 buenos; malos $=1.890$, el 38,18\% de los 4.950 malos; $\ln(64{,}29/38{,}18)\cdot100=+52{,}1$.

\textbf{La ley vista como fusión forzada.} Tras el borrón, los tres primeros atributos se ven iguales: quedan «sin registro visible» (94.000; impago 4,5\%; WOE $+10{,}0$) y «\$2,5 M o más» (WOE $-96{,}3$). El IV cae de \textbf{0,31 a 0,10}\calc: de fuerte a débil (Siddiqi: 0,02 / 0,1 / 0,3). Es exactamente lo que encuentra el paper: los morosos borrados se ven mejores (su riesgo predicho cae en promedio 32\%) y los que nunca debieron se ven peores (el 89\% sube, en promedio 29\%)\V.

\begin{enun}{T37 (núcleo) · signo del WOE}
V/F. El WOE de «mora de \$200 mil a \$2,5 M» es $-64{,}1$: en ese atributo hay proporcionalmente más buenos que malos.
\end{enun}
\begin{ejemplar}
\textbf{F.} WOE $<0\iff\%M>\%B$: el atributo reúne el 23,6\% de los malos y solo el 12,5\% de los buenos. Es un atributo de \textbf{más riesgo}. El único WOE positivo, $+52{,}1$ (sin registro), es el de menos riesgo.
\end{ejemplar}

\begin{enun}{T52 (núcleo) · con qué valor entra una persona}
Selección múltiple. En la regresión woeizada (antes de la ley), una persona con \$150 mil en mora entra con: a) 150.000 \quad b) dummy «hasta \$200 mil» $=1$ \quad c) $-0{,}289$ \quad d) $-28{,}9$
\end{enun}
\begin{ejemplar}
\textbf{c) $-0{,}289$.} Se la ubica en su atributo y entra con el \textbf{WOE de ese atributo, sin $\times100$}. No entra con el monto (eso sería otra especificación), ni con una dummy (regresión con dummies), ni con el WOE de la tabla ($\times100$). Después de la ley, esa misma persona entraría con $+0{,}100$.
\end{ejemplar}

\begin{enun}{T42 (refuerzo) · número de coeficientes}
V/F. Esta variable de 4 atributos aporta 4 coeficientes si entra como dummies y 1 si entra woeizada.
\end{enun}
\begin{ejemplar}
\textbf{F.} Como dummies aporta $k-1=\mathbf{3}$ (un atributo queda de base, en el intercepto). Woeizada aporta \textbf{1}: una sola columna cuyo valor cambia según el atributo.
\end{ejemplar}

\begin{enun}{T43 (refuerzo) · signo esperado de los $\hat\beta$ woeizados}
V/F. Con $y=1$ impago y el WOE de la clase, se espera que todos los $\hat\beta$ de la regresión woeizada sean positivos.
\end{enun}
\begin{ejemplar}
\textbf{F.} Se esperan \textbf{negativos}: el WOE crece con la calidad del deudor (más WOE, más «buenos») y el logit modela $y=1=$ malo, así que más WOE implica menos probabilidad de impago. Un $\hat\beta$ positivo es señal de multicolinealidad o de mala tramificación.
\end{ejemplar}

\begin{enun}{T59 (núcleo) · monotonía y fusión}
Comente. Para la carga financiera, \texttt{woebin} entrega 5 tramos (0--10, 10--20, 20--30, 30--43 y más de 43\% del ingreso) con tasas de impago de 2,5; 4,0; 3,5; 6,5 y 6,0\%\inv; $N_2=25.000$, $N_3=20.000$, $N_4=15.000$, $N_5=10.000$.
\end{enun}
\begin{ejemplar}
La tasa \textbf{no es monótona}: baja de 4,0 a 3,5\% y de 6,5 a 6,0\%, cuando a más carga se espera más impago. Se \textbf{retramifica fusionando tramos vecinos} y \textbf{ponderando por $N$}: tramo 2+3 $=\frac{25.000\cdot0,040+20.000\cdot0,035}{45.000}=3{,}78\%$; tramo 4+5 $=\frac{15.000\cdot0,065+10.000\cdot0,060}{25.000}=6{,}30\%$\calc. Queda $2{,}5\to3{,}78\to6{,}30$, monótona. La fusión de arriba tiene respaldo real: en el probit del Banco Central, sobre una carga de 42,8 la pendiente deja de ser significativa\V.
\end{ejemplar}
\begin{trampa}
\begin{itemize}[leftmargin=*,itemsep=0pt,topsep=0pt]
\item Nunca promediar tasas o WOE sin ponderar: $(4{,}0+3{,}5)/2=3{,}75$ está mal.
\item IV alto de una variable no es lo mismo que valor incremental: sin el registro, el predictor del paper pierde poco (log-verosimilitud $-2\%$) porque los bancos aún ven la mora bancaria y el historial de deuda\V.
\end{itemize}
\end{trampa}

\begin{enun}{T55 (refuerzo) · exclusiones}
Comente. Un banco excluye de su muestra de modelamiento a sus funcionarios y a quienes no tienen historial de crédito. Un diputado lo acusa de discriminar, «como hacía Dicom».
\end{enun}
\begin{ejemplar}
Las exclusiones se justifican por \textbf{calidad o representatividad}, no por características personales. Los funcionarios tienen condiciones preferenciales y otro incentivo a pagar: mezclarlos sesga la tasa de malos. Quien no tiene historial no aporta comportamiento de pago que observar; vuelve a la muestra cuando lo acumule. El paper hace lo mismo por diseño: deja fuera a quienes deben más de \$2,5 millones (su información no cambió) y exige saldo positivo 6 meses antes para tener una muestra comparable en el tiempo\V. Lo que la Ley 20.575 prohíbe es otra cosa: \emph{usar} el registro para fines ajenos al crédito, como contratar personal\Pre.
\end{ejemplar}

% =====================================================================
\acto{Acto 4. ¿Qué causó el borrón? Diferencias en diferencias}

\textbf{El diseño del paper}\V.
\begin{itemize}[leftmargin=*,itemsep=0pt,topsep=0pt]
\item \textbf{Exposición:} con un bosque aleatorio entrenado en dic.\ 2009 se predice el impago de cada persona con y sin el registro. $E_i=$ cambio porcentual de su riesgo predicho por el borrón.
\item \textbf{Tratados:} exposición positiva ($E_i\geq+15\%$; 2,1 millones, casi todos sin mora) y negativa ($E_i\leq-15\%$; 608 mil, casi todos morosos). \textbf{Control:} el grupo cero ($|E_i|<15\%$), a quienes el borrón no cambia la predicción.
\item \textbf{Antes/después:} cohortes de jun.\ 2010, dic.\ 2010 y jun.\ 2011 (placebos) y dic.\ 2011 (la tratada). $Y=$ crédito de consumo nuevo en los 6 meses siguientes (miles de \$). Errores agrupados en 330 «mercados».
\item \textbf{Resultado} (Tabla 4): exposición positiva, $\hat\beta_3=-13{,}72$ (EE 3,83), \textbf{$-6{,}4\%$} sobre una media de \$215 mil; exposición negativa, $+16{,}60$ (EE 3,72), \textbf{$+11{,}8\%$} sobre \$141 mil. Neto: $-\$30$ mil millones $+\ \$10$ mil millones $=-\$20$ mil millones en un semestre, \textbf{$-3{,}5\%$}\V.
\end{itemize}
Cuatro medias para practicar\inv, coherentes con la Tabla 4 (media previa del tratado y $\hat\beta_3$ reales) y con la Tabla 3 (grupo cero: \$175 mil en jul.\ 2011)\V:
\begin{center}\small
\begin{tabular}{lccc}\toprule
Crédito nuevo (miles de \$) & Antes (jun.\ 2011) & Después (dic.\ 2011) & Cambio\\\midrule
Exposición positiva (tratado) & 215,28 & 205,28 & $-10{,}00$\\
Grupo cero (control) & 175,00 & 178,72 & $+3{,}72$\\\midrule
DiD & & & $\mathbf{-13{,}72}\ \Rightarrow -13{,}72/215{,}28=-6{,}4\%$\\\bottomrule
\end{tabular}
\end{center}

\begin{enun}{T69 (núcleo) · qué es cada $\beta$}
V/F. En $y=\beta_0+\beta_1Treat+\beta_2Post+\beta_3(Treat\times Post)$, $\hat\beta_1=40{,}28$ es el efecto de la Ley Dicom en los que no tenían mora.
\end{enun}
\begin{ejemplar}
\textbf{F.} $\hat\beta_0=175{,}00$: control antes. $\hat\beta_1=215{,}28-175{,}00=40{,}28$: \textbf{brecha inicial} (los expuestos ya pedían más crédito). $\hat\beta_2=+3{,}72$: tendencia común. $\hat\beta_3=-13{,}72$: \textbf{el efecto} (\$13.720 menos por persona, $-6{,}4\%$). Elasticidad del crédito al riesgo predicho: $-0{,}064/0{,}22=-0{,}29$\V\calc.
\end{ejemplar}

\begin{enun}{T62 (núcleo) · niveles distintos no sesgan}
V/F. Los expuestos negativos piden \$141 mil y los positivos \$215 mil antes de la ley; como los niveles son tan distintos del grupo cero, el DiD está sesgado.
\end{enun}
\begin{ejemplar}
\textbf{F} (y T62, tal como está en el banco, es \textbf{V}). Tendencias paralelas pide que, \textbf{sin la ley}, ambos grupos hubieran \textbf{cambiado} lo mismo, no que tengan el mismo nivel. La brecha permanente ($\beta_1$) se elimina al diferenciar. Lo que sesga son tendencias distintas.
\end{ejemplar}

\begin{enun}{T63 (refuerzo) · cómo se revisa el supuesto}
V/F. Si el paper tuviera solo jun.\ 2011 y dic.\ 2011, verificaría tendencias paralelas comparando el crédito de ambos grupos en jun.\ 2011.
\end{enun}
\begin{ejemplar}
\textbf{F.} Con dos periodos \textbf{no se puede verificar}: es un contrafactual, y comparar niveles no dice nada de tendencias. El paper tiene dos cohortes previas más y mira las \textbf{tendencias previas}: con jun.\ 2011 $=0$, los coeficientes de jun.\ y dic.\ 2010 son $-4{,}67$ (2,81) y $-0{,}25$ (3,25) en exposición positiva, y $-4{,}74$ (2,30) y $0{,}75$ (2,59) en la negativa\V. Ojo crítico: $-4{,}74/2{,}30$ da $t=-2{,}06$\calc, significativo al 5\%; los autores los describen como planos. Por eso agregan dos pruebas: un DiD alrededor del corte de \$2,5 millones y una comparación con los mismos meses de años anteriores\V.
\end{ejemplar}

\begin{enun}{T61 (refuerzo) · ATT frente a ATE}
V/F. En un experimento aleatorio la diferencia de medias estima el ATT, y el DiD de la Ley Dicom estima el efecto promedio en toda la población (ATE).
\end{enun}
\begin{ejemplar}
\textbf{F.} Es al revés. El experimento aleatorio da el \textbf{ATE}; el DiD estima el \textbf{ATT}: el efecto en los \textbf{tratados} (los 2,1 millones con exposición positiva y los 608 mil con exposición negativa), no en un deudor cualquiera. Los expuestos negativos, por ejemplo, son más pobres que el promedio (71\% de nivel socioeconómico bajo contra 55\%)\V: su $+11{,}8\%$ no se traslada a toda la población.
\end{ejemplar}

\begin{enun}{T66 (refuerzo) · interna frente a externa}
V/F. Los efectos de equilibrio general y las muestras pequeñas (p.\,ej., usar solo a quienes deben entre \$2,25 y \$2,75 millones) son amenazas a la validez externa.
\end{enun}
\begin{ejemplar}
\textbf{F (a medias).} El \textbf{equilibrio general} es amenaza \textbf{externa}: un borrón chico, para unos pocos, no movería los precios; uno para el 21\% de los adultos hace que los bancos re-predigan el riesgo de todos y suban las tasas (el corazón del paper). También es externa la especificidad del caso: Chile 2012, donde los bancos aún veían la mora bancaria. La \textbf{muestra pequeña} amenaza la validez \textbf{interna}: con ventanas muy estrechas alrededor del corte «los errores estándar crecen», reconocen los autores\V.
\end{ejemplar}

\textbf{Las amenazas, con el caso.} \textbf{Derrame a los controles:} el borrón afectó también a quienes no tenían mora (su riesgo predicho subió); por eso el control no puede ser «los no borrados», sino el grupo cero, y se revisa con el corte de \$2,5 M\V. \textbf{Shock simultáneo:} si algo más hubiera cambiado en 2012, también habrían subido las tasas de los créditos grandes, y no subieron\V; el terremoto de 2010 no elevó la mora agregada\V. \textbf{Composición:} se exige saldo positivo 6 meses antes, para que no entren solo quienes después pidieron crédito\V. \textbf{Anticipación:} la ley se discutió en 2011 con fecha de corte conocida; habría que revisar si alguien dejó de pagar para caer bajo el borrón (amenaza por revisar, no un hallazgo del paper).

\begin{enun}{T77 (núcleo) · la nota de prensa falla y se propone el DiD}
Comente. Una nota dice: «Tras la Ley Dicom, los beneficiados volvieron a endeudarse: un año después, piden más créditos que antes». Otra compara, \emph{después} de la ley, el crédito de los beneficiados con el de quienes nunca estuvieron en Dicom. Ambas concluyen que «la ley amplió el acceso al crédito». Comente y proponga cómo evaluarlo.
\end{enun}
\begin{ejemplar}
\textbf{La primera es un antes/después sin control}: atribuye a la ley todo lo que cambió en un año (ciclo, tasas, estacionalidad). \textbf{La segunda es un después sin antes}: los beneficiados ya eran distintos (más pobres, con menos historial), así que la brecha mezcla el efecto con diferencias previas (\textbf{selección}). Y ninguna mira a los perjudicados. Lo correcto es un \textbf{DiD}: el cambio del crédito de los expuestos, menos el cambio de un grupo comparable \emph{no afectado} (el grupo cero), con varias cohortes previas. \textbf{Supuesto:} tendencias paralelas, que se revisa con las tendencias previas. \textbf{Amenazas:} derrame (los «no beneficiados» también fueron afectados), shock simultáneo y anticipación. El resultado contradice a la prensa: los beneficiados ganan $+11{,}8\%$, pero los perjudicados son más y el total cae $-3{,}5\%$\V. \textbf{En la PEP, con el SLEP, el esqueleto es el mismo:} «después sin antes» → DiD con colegios que siguen municipales como control, cambio del SIMCE antes y después del traspaso, tendencias previas del SIMCE y amenazas (selección de qué comunas se traspasan primero, shocks simultáneos).
\end{ejemplar}

\textbf{Mini ejercicio: el DiD ingenuo de la Figura 4.} El paper parte comparando morosos (borrados) con no morosos: el crédito de los morosos sube \textbf{\$41 mil} más que el de los no morosos, un 46\% de su media previa de \$88 mil\V\calc. Con medias inventadas para los no morosos\inv: morosos $88\to119$ ($+31$); no morosos $200\to190$ ($-10$); DiD $=31-(-10)=\mathbf{+41}$. Dos lecturas de examen. (1) El control \emph{también fue tratado}: sus $-10$ son en parte efecto de la ley (su riesgo predicho subió). Si sin la ley se hubieran quedado en 200, el efecto real en los morosos sería $+31$: el DiD ingenuo \textbf{sobrestima} la ganancia. (2) Este DiD compara dos grupos afectados en direcciones opuestas, así que no dice nada del efecto total. Por eso los autores pasan al grupo cero y llegan a $-3{,}5\%$; en su modelo de selección adversa, la pérdida de excedente sube un 66\% respecto de la situación previa\V.

% =====================================================================
\newpage
\acto{Tabla de recuerdo activo: tapa la columna derecha}
\small
\renewcommand{\arraystretch}{1.18}
\begin{tabular}{@{}p{1.1cm}p{6.6cm}p{8.7cm}@{}}\toprule
Ítem & En el caso & Respuesta en una línea (palabra clave en negrita)\\\midrule
T14 & Probit BCCh: carga financiera con $\hat\beta>0$ & F · $\hat\beta>0$ con $y=1$ malo $=$ \textbf{más riesgo}\\
T6 & $\hat\beta_{REG}=0{,}909$ «son 0,91 pp» & F · \textbf{odds ×2,48}; efecto dummy $=\Ph(1)-\Ph(0)=+4{,}0$ pp\\
T13 & Tasa 26 $\to$ 31,3\% «sube 5,3\%» & F · \textbf{5,3 pp}; relativo $+20\%$\\
T1 & MPL, varianza 2,3 veces mayor & F · \textbf{insesgado}, pierde eficiencia; errores robustos\\
T22 & Despido: $\hat\beta$ igual en $P=0{,}1$ y $0{,}5$ & F · $g(z)\hat\beta$: $+1{,}4$ contra $+3{,}2$ pp; \textbf{el del margen}\\
T11 & 4 deudores, EPeP $=$ EPP & F · $+0{,}136$ contra $+0{,}165$ pp; \textbf{$g$ no lineal}\\
T33 & Matriz 40/460/20/9.480 & Exact.\ 95,2 · sens.\ 8 · espec.\ 99,8 · prec.\ 66,7 · F1 14,3; sobra $1-$prec.\\
T34 & «Acierta el 95,2\%» & \textbf{Tasa base} 95\% (Fed: 85\%); mirar sensibilidad, ROC, AUC\\
T23 & Corte 0,5 «el mejor» & F · depende de \textbf{costos}; con $\bar y=0{,}05$: sens.\ 70\%\\
T27 & ROC «S contra E» & F · S contra \textbf{1 $-$ E}; la ley \textbf{baja la curva}\\
T37 & WOE $-64$ «más buenos» & F · WOE $<0$: \textbf{\%M $>$ \%B}, más riesgo\\
T52 & \$150 mil en mora entra con... & \textbf{$-0{,}289$} (WOE sin ×100)\\
T42 & 4 atributos: ¿4 y 1 $\beta$? & F · dummies $k-1=\mathbf{3}$; woeizada 1\\
T43 & $\hat\beta$ woeizados positivos & F · se esperan \textbf{negativos}\\
T59 & 2,5; 4,0; 3,5; 6,5; 6,0\% & No monótona: \textbf{fusionar ponderando por $N$}: 2,5; 3,78; 6,30\\
T55 & Funcionarios y sin historial & \textbf{Calidad o representatividad}, no discriminación\\
T69 & $\hat\beta_1=40{,}28$ «es el efecto» & F · $\beta_1$ \textbf{brecha inicial}; efecto $=\beta_3=-13{,}72$ ($-6{,}4\%$)\\
T62 & 141 y 215 frente a 175: «sesgado» & F (el T62 del banco es V) · niveles se van al diferenciar; sesgan las \textbf{tendencias}\\
T63 & 2 periodos, comparar niveles & F · \textbf{no se puede verificar}; con más cohortes, tendencias previas\\
T61 & ECA $=$ ATT, DiD $=$ ATE & F · al revés: ECA $=$ ATE, DiD $=$ \textbf{ATT}\\
T66 & Equilibrio general y muestra chica, externas & F a medias · muestra pequeña es \textbf{interna}; equilibrio general es externa\\
T77 & «La ley amplió el crédito» & Antes/después y después solo: \textbf{selección}; \textbf{DiD} + tendencias previas + amenazas; SLEP igual\\\bottomrule
\end{tabular}

\normalsize
\vspace{8pt}
\textbf{Fuentes}
\footnotesize\sloppy
\begin{itemize}[leftmargin=*,itemsep=0pt,topsep=0pt]
\item Liberman, A., Neilson, C., Opazo, L. y Zimmerman, S. (2018, rev.\ 2019). The Equilibrium Effects of Information Deletion: Evidence from Consumer Credit Markets. NBER Working Paper 25097 (secciones 2--4; tablas 1--4, 6 y 7; versión de octubre de 2019 en la web de Neilson, mismas cifras). No publicado en revista: la web del autor lo lista como \emph{revise and resubmit} en el \emph{Journal of Finance} y resume una versión más nueva con otras cifras (crédito agregado $-9\%$, elasticidad $-0{,}14$, 67\% de los morosos), cuyo texto no se pudo revisar. La ficha usa la versión de 2019.
\item Martínez, F., Cifuentes, R., Madeira, C. y Poblete-Cazenave, R. (2013). Measurement of Household Financial Risk with the Survey of Household Finances. Documento de Trabajo 682, Banco Central de Chile (tabla 5).
\item Ley 20.575 (2012), que establece el principio de finalidad en el tratamiento de datos personales. Detalle de la ley: Cooperativa.cl, 24-1-2012.
\item Estrella, A. y Mishkin, F. (1998). Predicting U.S. Recessions: Financial Variables as Leading Indicators. \emph{Review of Economics and Statistics}, 80(1), 45--61. Coeficientes del modelo de la Fed de Nueva York tomados de fuentes secundarias\Pre; muestra 1959--2009 y probabilidad de julio de 2023 en Neely, C.\ (2023), «What Is the Probability of a Recession?», \emph{On the Economy}, Fed de St.\ Louis, 7-9-2023.
\item NBER, US Business Cycle Expansions and Contractions (fechas de recesiones; base del cálculo de 93 de 612 meses).
\item Kaminsky, G. y Reinhart, C. (1999). The Twin Crises: The Causes of Banking and Balance-of-Payments Problems. \emph{American Economic Review}, 89(3), 473--500.
\item Siddiqi, N. (2006). \emph{Credit Risk Scorecards}. Wiley (umbrales del IV).
\end{itemize}
\end{document}
